% !TEX program = xelatex
% Compile with XeLaTeX for exact Times New Roman / Courier New fonts.
% The pdfLaTeX fallback prevents fontspec errors in default LaTeX editors.
% This file intentionally has no title page.
% Self-contained report: source code is embedded, so no external C/H files are required for the Program Code section.

\documentclass[12pt,a4paper]{article}

% --- Page + fonts ---
\usepackage[a4paper,margin=2cm]{geometry}
\usepackage{iftex}
\usepackage{setspace}
\ifPDFTeX
    \usepackage[T1]{fontenc}
    \usepackage[utf8]{inputenc}
    \usepackage{mathptmx}
    \usepackage{courier}
\else
    \usepackage{fontspec}
    \setmainfont{Times New Roman}
    \setmonofont{Courier New}
\fi

% --- Header/footer (page number bottom-center) ---
\usepackage{fancyhdr}
\pagestyle{fancy}
\fancyhf{}
\fancyfoot[C]{\thepage}
\renewcommand{\headrulewidth}{0pt}

% --- Lists, tables ---
\usepackage{enumitem}
\usepackage{longtable}
\usepackage{array}

% --- Images + float control ---
\usepackage{graphicx}
\usepackage{float}
\usepackage{placeins}
\usepackage{caption}
\DeclareCaptionFont{t14}{\fontsize{14}{21}\selectfont}
\DeclareCaptionFont{t12}{\fontsize{12}{18}\selectfont}
\captionsetup{font=t14}

% --- Boxes ---
\usepackage[most]{tcolorbox}
\tcbuselibrary{breakable}

% --- Code formatting (VS Code-like, light) ---
\usepackage{xcolor}
\usepackage{listings}

\definecolor{vscKeyword}{RGB}{0,0,180}
\definecolor{vscComment}{RGB}{0,128,0}
\definecolor{vscString}{RGB}{163,21,21}
\definecolor{vscNumber}{RGB}{9,134,88}
\definecolor{softGray}{RGB}{245,245,245}

\lstdefinestyle{vscodeC}{
  language=C,
  basicstyle=\ttfamily\fontsize{9.5}{11}\selectfont\linespread{1}\selectfont,
  columns=fullflexible,
  breaklines=true,
  showstringspaces=false,
  tabsize=4,
  numbers=left,
  numberstyle=\ttfamily\fontsize{8}{9}\selectfont\color{gray},
  frame=single,
  rulecolor=\color{black},
  backgroundcolor=\color{softGray},
  keywordstyle=\color{vscKeyword}\bfseries,
  commentstyle=\color{vscComment},
  stringstyle=\color{vscString},
  literate=%
   *{0}{{{\color{vscNumber}0}}}{1}
    {1}{{{\color{vscNumber}1}}}{1}
    {2}{{{\color{vscNumber}2}}}{1}
    {3}{{{\color{vscNumber}3}}}{1}
    {4}{{{\color{vscNumber}4}}}{1}
    {5}{{{\color{vscNumber}5}}}{1}
    {6}{{{\color{vscNumber}6}}}{1}
    {7}{{{\color{vscNumber}7}}}{1}
    {8}{{{\color{vscNumber}8}}}{1}
    {9}{{{\color{vscNumber}9}}}{1}
}

\newcommand{\reportimage}[3]{
\begin{figure}[H]
    \centering
    \IfFileExists{#1}{
        \includegraphics[width=#2\linewidth]{#1}
    }{
        \fbox{\parbox[c][5cm][c]{0.88\linewidth}{\centering Missing screenshot file: \texttt{\detokenize{#1}}}}
    }
    \caption{#3}
\end{figure}
}

\begin{document}

\setcounter{page}{1}

{\fontsize{14}{21}\selectfont
\setstretch{1.5}

\tableofcontents
\clearpage

\section{Task}

\noindent Solve the given laboratory problem in C by using user-defined data types, pointers, dynamic memory
allocation, Stack and Queue abstract data types, and the \texttt{FILE} data type. The solution must be written
in a procedural style, using self-written functions and header files, and it must be executable through a
complete menu from the \texttt{main} function.

\vspace{0.45cm}
\noindent\textbf{Problem Condition:}

\begin{center}
\begin{tcolorbox}[
  breakable,
  width=0.90\textwidth,
  colback=white,
  colframe=black,
  boxrule=0.8pt,
  left=10pt,right=10pt,top=10pt,bottom=10pt
]
\setstretch{1.5}

Develop a procedural-style program in C/C++ using own written functions. Data processing in the program must be
organized according to a given length of input records based on memory allocation functions.

The program must use pointers, custom data types based on structures, nested structures where necessary, and
pointer-specific notation. The solution must be organized through header files and must run through a complete
menu in the \texttt{main} function.

Two distinct versions must be presented:

\begin{enumerate}[label=\Alph*., leftmargin=1.2cm, itemsep=0.25em]
    \item A version implemented as a dynamic Stack based on a List ADT.
    \item A version implemented as dynamic Queues based on a List ADT, including Simple Queue, Double Ended
    Queue, Circular Queue, and Priority Queue.
\end{enumerate}

The program must support creating, traversing and displaying, inserting, searching by position or value,
deleting, and registering the newly formed Stack/Queue into files in both text and binary mode. Files must be
possible to create, open, reopen, and delete.

\end{tcolorbox}
\end{center}

\begin{center}
\textit{Figure 1.1 (Problem statement)}
\end{center}

\clearpage

\section{Proposed Solution}

The developed program uses a custom data type named \texttt{Citizen}. This type represents the input and output
record processed by all Stack and Queue operations. Each citizen contains personal data, a date of birth, two
nested address structures, and computed output fields such as age, category, and amount paid.

The solution is divided into modules:

\begin{longtable}{|p{0.24\textwidth}|p{0.66\textwidth}|}
\hline
\textbf{File} & \textbf{Responsibility} \\
\hline
\texttt{data\_types.h} and \texttt{data\_types.c} & Define the user-defined data types and display helpers. \\
\hline
\texttt{stack.h} and \texttt{stack.c} & Implement Version A, the dynamic Stack based on a linked list. \\
\hline
\texttt{queue.h} and \texttt{queue.c} & Implement Version B, the Simple Queue, Double Ended Queue, Circular Queue, and Priority Queue. \\
\hline
\texttt{file\_io.h} and \texttt{file\_io.c} & Implement file operations for text files, binary files, loading, reopening, and deleting. \\
\hline
\texttt{main.c} & Implements input validation and the complete menu system that calls the modules. \\
\hline
\end{longtable}

The implementation follows a procedural style: all operations are represented as functions, and the program
state is passed through pointers to structures such as \texttt{Stack*}, \texttt{Queue*}, and \texttt{Citizen*}.

\subsection{User-Defined Data Type}

The central user-defined data type is:

\begin{lstlisting}[style=vscodeC]
typedef struct {
    int day;
    int month;
    int year;
} Date;

typedef struct {
    char city[CITY_LENGTH];
    char street[STREET_LENGTH];
    char postCode[POST_CODE_LENGTH];
} Address;

typedef struct {
    char name[NAME_LENGTH];
    char surname[SURNAME_LENGTH];
    Date dob;
    char gender;
    Address home;
    Address work;

    int ageYears;
    int ageMonths;
    int ageDays;
    int category;
    float amountPaid;
} Citizen;
\end{lstlisting}

The structure \texttt{Citizen} contains nested structures: \texttt{Date}, \texttt{home Address}, and
\texttt{work Address}. This satisfies the requirement to use structures and nested structures.

\subsection{Linked List Node}

The Stack and Queue implementations use dynamically allocated linked-list nodes:

\begin{lstlisting}[style=vscodeC]
typedef struct Node {
    Citizen data;
    int priority;
    struct Node* next;
    struct Node* prev;
} Node;
\end{lstlisting}

The fields \texttt{next} and \texttt{prev} are pointers. They allow the same node structure to support Stack,
Simple Queue, Deque, Circular Queue, and Priority Queue behavior.

\clearpage

\section{Version A: Dynamic Stack Based on List ADT}

Version A is implemented through the \texttt{Stack} structure:

\begin{lstlisting}[style=vscodeC]
typedef struct {
    Node* top;
    size_t size;
} Stack;
\end{lstlisting}

The Stack stores a pointer to the top node and the current number of records. Every inserted element is
allocated using \texttt{malloc}. Every deleted element is released using \texttt{free}.

\subsection{Operations Implemented}

\begin{longtable}{|p{0.33\textwidth}|p{0.57\textwidth}|}
\hline
\textbf{Required operation} & \textbf{Implemented function(s)} \\
\hline
Create Stack & \texttt{initStack}, \texttt{createStackFromRecords} \\
\hline
Traverse/display Stack & \texttt{displayStack} \\
\hline
Insert element & \texttt{pushStack}, called from the menu as push operation \\
\hline
Search by position & \texttt{searchStackByPosition} \\
\hline
Search by value & \texttt{searchStackBySurname} \\
\hline
Delete element & \texttt{popStack}, \texttt{deleteStackByPosition}, \texttt{deleteStackBySurname} \\
\hline
Register in file & \texttt{saveStackToTextFile}, \texttt{saveStackToBinaryFile} \\
\hline
Reopen/load from file & \texttt{loadStackFromBinaryFile}, \texttt{displayStackBinaryFile} \\
\hline
\end{longtable}

\subsection{Stack Insertion}

\begin{lstlisting}[style=vscodeC]
int pushStack(Stack* s, const Citizen* data) {
    Node* node;

    if (s == NULL || data == NULL) {
        return 0;
    }

    node = createNode(data, 0);
    if (node == NULL) {
        printf("Memory allocation failed. Stack insert was cancelled.\n");
        return 0;
    }

    node->next = s->top;
    if (s->top != NULL) {
        s->top->prev = node;
    }

    s->top = node;
    s->size++;
    return 1;
}
\end{lstlisting}

This function inserts a new record at the top of the Stack. It uses pointer notation, dynamic allocation, and
linked-list references.

\clearpage

\section{Version B: Dynamic Queues Based on List ADT}

Version B is implemented through one unified \texttt{Queue} structure:

\begin{lstlisting}[style=vscodeC]
typedef enum {
    QUEUE_SIMPLE = 1,
    QUEUE_DEQUE = 2,
    QUEUE_CIRCULAR = 3,
    QUEUE_PRIORITY = 4
} QueueType;

typedef struct {
    Node* front;
    Node* rear;
    size_t size;
    QueueType type;
} Queue;
\end{lstlisting}

The same structure supports all four required queue types. The selected queue behavior is controlled by the
\texttt{type} field.

\subsection{Queue Types Implemented}

\begin{longtable}{|p{0.28\textwidth}|p{0.62\textwidth}|}
\hline
\textbf{Queue type} & \textbf{Behavior} \\
\hline
Simple Queue & Inserts at the rear and deletes from the front. \\
\hline
Double Ended Queue & Supports insertion at front/rear and deletion from front/rear. \\
\hline
Circular Queue & The rear node points back to the front node, and the front node points back to the rear. \\
\hline
Priority Queue & Records are kept ordered by priority; priority \texttt{1} is the highest priority. \\
\hline
\end{longtable}

\subsection{Operations Implemented}

\begin{longtable}{|p{0.33\textwidth}|p{0.57\textwidth}|}
\hline
\textbf{Required operation} & \textbf{Implemented function(s)} \\
\hline
Create Queue & \texttt{initQueue}, \texttt{createQueueFromRecords} \\
\hline
Traverse/display Queue & \texttt{displayQueue} \\
\hline
Insert element & \texttt{enqueueQueue}, \texttt{enqueueQueueFront} \\
\hline
Search by position & \texttt{searchQueueByPosition} \\
\hline
Search by value & \texttt{searchQueueBySurname} \\
\hline
Delete element & \texttt{dequeueQueue}, \texttt{dequeueQueueRear}, \texttt{deleteQueueByPosition}, \texttt{deleteQueueBySurname} \\
\hline
Register in file & \texttt{saveQueueToTextFile}, \texttt{saveQueueToBinaryFile} \\
\hline
Reopen/load from file & \texttt{loadQueueFromBinaryFile}, \texttt{displayQueueBinaryFile} \\
\hline
\end{longtable}

\subsection{Priority Queue Insertion}

\begin{lstlisting}[style=vscodeC]
static int enqueueByPriority(Queue* q, const Citizen* data, int priority) {
    Node* node;
    Node* current;

    node = createNode(data, priority);
    if (node == NULL) {
        printf("Memory allocation failed. Priority queue insert was cancelled.\n");
        return 0;
    }

    if (isQueueEmpty(q)) {
        q->front = node;
        q->rear = node;
    } else if (priority < q->front->priority) {
        node->next = q->front;
        q->front->prev = node;
        q->front = node;
    } else {
        current = q->front;
        while (current->next != NULL && current->next->priority <= priority) {
            current = current->next;
        }

        node->next = current->next;
        node->prev = current;
        if (current->next != NULL) {
            current->next->prev = node;
        } else {
            q->rear = node;
        }
        current->next = node;
    }

    q->size++;
    return 1;
}
\end{lstlisting}

The Priority Queue inserts each new node into the correct position according to its priority. This makes the
front node always represent the highest-priority record.

\clearpage

\section{File Processing}

The file processing module uses the C \texttt{FILE} data type and standard file operations:
\texttt{fopen}, \texttt{fclose}, \texttt{fprintf}, \texttt{fwrite}, \texttt{fread}, \texttt{fgetc}, and
\texttt{remove}.

\subsection{Supported File Operations}

\begin{enumerate}[label=(\roman*), leftmargin=1.2cm, itemsep=0.25em]
    \item Create an empty text file.
    \item Create an empty binary file.
    \item Open or reopen a text file and display its contents.
    \item Open or reopen a binary Stack file.
    \item Open or reopen a binary Queue file.
    \item Delete a selected file.
    \item Save Stack records in text mode.
    \item Save Stack records in binary mode.
    \item Save Queue records in text mode.
    \item Save Queue records in binary mode.
\end{enumerate}

\subsection{Example of Saving a Queue in Binary Mode}

\begin{lstlisting}[style=vscodeC]
int saveQueueToBinaryFile(const Queue* q, const char* filename) {
    FILE* file;
    Node* current;
    QueueFileRecord record;
    size_t position;

    if (q == NULL || filename == NULL) {
        return 0;
    }

    file = fopen(filename, "wb");
    if (file == NULL) {
        printf("Could not open file for binary write: %s\n", filename);
        return 0;
    }

    current = q->front;
    for (position = 1; position <= q->size; position++) {
        record.data = current->data;
        record.priority = current->priority;
        fwrite(&record, sizeof(QueueFileRecord), 1, file);
        current = current->next;
    }

    fclose(file);
    printf("Queue saved in binary mode: %s\n", filename);
    return 1;
}
\end{lstlisting}

\clearpage

\section{Requirement Compliance}

\begin{longtable}{|p{0.07\textwidth}|p{0.48\textwidth}|p{0.35\textwidth}|}
\hline
\textbf{No.} & \textbf{Requirement} & \textbf{Status} \\
\hline
1 & Use pointers. & Implemented through \texttt{Node*}, \texttt{Stack*}, \texttt{Queue*}, \texttt{Citizen*}. \\
\hline
2 & Use structures and user-defined data types. & Implemented with \texttt{Citizen}, \texttt{Date}, \texttt{Address}, \texttt{Node}, \texttt{Stack}, \texttt{Queue}. \\
\hline
3 & Use nested structures where necessary. & \texttt{Citizen} contains \texttt{Date}, \texttt{home Address}, and \texttt{work Address}. \\
\hline
4 & Use dynamic memory allocation. & Implemented with \texttt{malloc}, \texttt{realloc}, and \texttt{free}. \\
\hline
5 & Organize solution through header files. & Implemented with \texttt{data\_types.h}, \texttt{stack.h}, \texttt{queue.h}, and \texttt{file\_io.h}. \\
\hline
6 & Run through a complete menu in \texttt{main}. & Implemented in \texttt{main.c}. \\
\hline
7 & Version A: dynamic Stack based on List ADT. & Implemented in \texttt{stack.c}. \\
\hline
8 & Version B: dynamic Queues based on List ADT. & Implemented in \texttt{queue.c}. \\
\hline
9 & Include Simple, Double Ended, Circular, and Priority Queue. & All four types are available through the Version B menu. \\
\hline
10 & Create, traverse, and display. & Implemented for Stack and all Queues. \\
\hline
11 & Insert elements. & Implemented for Stack and all Queues. \\
\hline
12 & Search by position or value. & Implemented by position and by surname. \\
\hline
13 & Delete elements. & Implemented by ADT operation, by position, and by surname. \\
\hline
14 & Register records in text and binary files. & Implemented through \texttt{file\_io.c}. \\
\hline
15 & Create, open, reopen, and delete files. & Implemented through the File Tools menu. \\
\hline
16 & Commit to Git/GitHub. & Completed. Last commit: \texttt{76cdff9 Implement stack and queue lab versions}. \\
\hline
\end{longtable}

\clearpage

\section{Block Diagram}

The block diagram corresponding to the solved problem should show the main menu, the two versions, the four
queue types, the file tools, and the exit path with memory release.

\reportimage{block_diagram.png}{0.95}{Block diagram of the Stack/Queue program.}

\clearpage

\section{Program Outputs}
\captionsetup{font=t12}

\reportimage{fig_01_main_menu.png}{0.95}{Main menu with Version A, Version B, file tools, and exit option.}

\reportimage{fig_02_version_a_menu.png}{0.95}{Version A menu for the dynamic Stack based on List ADT.}

\reportimage{fig_03_stack_create_display.png}{0.95}{Creating records and displaying the dynamic Stack.}

\reportimage{fig_04_stack_search_delete.png}{0.95}{Searching and deleting a Stack element by position or surname.}

\reportimage{fig_05_stack_text_binary_save.png}{0.95}{Saving the Stack in text mode and binary mode.}

\reportimage{fig_06_version_b_queue_type.png}{0.95}{Selecting one of the four Queue types in Version B.}

\reportimage{fig_07_simple_queue.png}{0.95}{Simple Queue insertion, deletion, and traversal.}

\reportimage{fig_08_deque.png}{0.95}{Double Ended Queue operations: insertion/deletion from front and rear.}

\reportimage{fig_09_circular_queue.png}{0.95}{Circular Queue traversal after insertion and deletion.}

\reportimage{fig_10_priority_queue.png}{0.95}{Priority Queue ordering, where priority 1 is processed first.}

\reportimage{fig_11_queue_files.png}{0.95}{Registering Queue data into text and binary files.}

\reportimage{fig_12_file_tools.png}{0.95}{File tools: create, open/reopen, and delete files.}

\captionsetup{font=t14}

\clearpage

\section{Program Code}

\noindent\textbf{Compilation command:}

\begin{lstlisting}[style=vscodeC]
gcc -std=c11 -Wall -Wextra -pedantic main.c data_types.c stack.c queue.c file_io.c -o program.exe
\end{lstlisting}

\subsection{data\_types.h}
\begin{lstlisting}[style=vscodeC]
#ifndef DATA_TYPES_H
#define DATA_TYPES_H

#include <stddef.h>

#define NAME_LENGTH 50
#define SURNAME_LENGTH 50
#define CITY_LENGTH 50
#define STREET_LENGTH 50
#define POST_CODE_LENGTH 20

typedef enum {
    CATEGORY_CHILD = 1,
    CATEGORY_ADULT = 2,
    CATEGORY_SENIOR = 3
} CitizenCategory;

typedef struct {
    int day;
    int month;
    int year;
} Date;

typedef struct {
    char city[CITY_LENGTH];
    char street[STREET_LENGTH];
    char postCode[POST_CODE_LENGTH];
} Address;

typedef struct {
    char name[NAME_LENGTH];
    char surname[SURNAME_LENGTH];
    Date dob;
    char gender;
    Address home;
    Address work;

    int ageYears;
    int ageMonths;
    int ageDays;
    int category;
    float amountPaid;
} Citizen;

typedef struct Node {
    Citizen data;
    int priority;
    struct Node* next;
    struct Node* prev;
} Node;

const char* citizenCategoryName(int category);
void printCitizenBrief(const Citizen* citizen, size_t position, int priority, int showPriority);
void printCitizenDetails(const Citizen* citizen);

#endif

\end{lstlisting}

\subsection{data\_types.c}
\begin{lstlisting}[style=vscodeC]
#include <stdio.h>
#include "data_types.h"

const char* citizenCategoryName(int category) {
    switch (category) {
        case CATEGORY_CHILD:
            return "Child";
        case CATEGORY_ADULT:
            return "Adult";
        case CATEGORY_SENIOR:
            return "Senior";
        default:
            return "Unknown";
    }
}

void printCitizenBrief(const Citizen* citizen, size_t position, int priority, int showPriority) {
    if (citizen == NULL) {
        return;
    }

    printf("%3zu. %-16s %-16s | DOB: %02d/%02d/%04d | Age: %02dY %02dM %02dD | %s",
           position,
           citizen->name,
           citizen->surname,
           citizen->dob.day,
           citizen->dob.month,
           citizen->dob.year,
           citizen->ageYears,
           citizen->ageMonths,
           citizen->ageDays,
           citizenCategoryName(citizen->category));

    if (showPriority) {
        printf(" | Priority: %d", priority);
    }

    printf("\n");
}

void printCitizenDetails(const Citizen* citizen) {
    if (citizen == NULL) {
        return;
    }

    printf("\nCitizen record\n");
    printf("Name: %s %s\n", citizen->name, citizen->surname);
    printf("Date of birth: %02d/%02d/%04d\n",
           citizen->dob.day,
           citizen->dob.month,
           citizen->dob.year);
    printf("Gender: %c\n", citizen->gender);
    printf("Home address: %s, %s, %s\n",
           citizen->home.city,
           citizen->home.street,
           citizen->home.postCode);
    printf("Work address: %s, %s, %s\n",
           citizen->work.city,
           citizen->work.street,
           citizen->work.postCode);
    printf("Computed age: %d years, %d months, %d days\n",
           citizen->ageYears,
           citizen->ageMonths,
           citizen->ageDays);
    printf("Category: %s\n", citizenCategoryName(citizen->category));
    printf("Amount paid: %.2f\n", citizen->amountPaid);
}

\end{lstlisting}

\subsection{stack.h}
\begin{lstlisting}[style=vscodeC]
#ifndef STACK_H
#define STACK_H

#include "data_types.h"

typedef struct {
    Node* top;
    size_t size;
} Stack;

void initStack(Stack* s);
int isStackEmpty(const Stack* s);
size_t stackSize(const Stack* s);

int pushStack(Stack* s, const Citizen* data);
int popStack(Stack* s, Citizen* removed);

int searchStackByPosition(const Stack* s, size_t position, Citizen* found);
int searchStackBySurname(const Stack* s, const char* surname, Citizen* found, size_t* position);

int deleteStackByPosition(Stack* s, size_t position, Citizen* removed);
int deleteStackBySurname(Stack* s, const char* surname, Citizen* removed, size_t* position);

void displayStack(const Stack* s);
void clearStack(Stack* s);

#endif

\end{lstlisting}

\subsection{stack.c}
\begin{lstlisting}[style=vscodeC]
#include <stdio.h>
#include <stdlib.h>
#include <string.h>
#include "stack.h"

static Node* createNode(const Citizen* data, int priority) {
    Node* node = (Node*)malloc(sizeof(Node));
    if (node == NULL) {
        return NULL;
    }

    node->data = *data;
    node->priority = priority;
    node->next = NULL;
    node->prev = NULL;
    return node;
}

static Node* nodeAtPosition(const Stack* s, size_t position) {
    Node* current;
    size_t index;

    if (s == NULL || position == 0 || position > s->size) {
        return NULL;
    }

    current = s->top;
    for (index = 1; index < position; index++) {
        current = current->next;
    }

    return current;
}

void initStack(Stack* s) {
    if (s == NULL) {
        return;
    }

    s->top = NULL;
    s->size = 0;
}

int isStackEmpty(const Stack* s) {
    return s == NULL || s->top == NULL;
}

size_t stackSize(const Stack* s) {
    return s == NULL ? 0 : s->size;
}

int pushStack(Stack* s, const Citizen* data) {
    Node* node;

    if (s == NULL || data == NULL) {
        return 0;
    }

    node = createNode(data, 0);
    if (node == NULL) {
        printf("Memory allocation failed. Stack insert was cancelled.\n");
        return 0;
    }

    node->next = s->top;
    if (s->top != NULL) {
        s->top->prev = node;
    }

    s->top = node;
    s->size++;
    return 1;
}

int popStack(Stack* s, Citizen* removed) {
    Node* temp;

    if (isStackEmpty(s)) {
        printf("Stack underflow. The stack is empty.\n");
        return 0;
    }

    temp = s->top;
    if (removed != NULL) {
        *removed = temp->data;
    }

    s->top = temp->next;
    if (s->top != NULL) {
        s->top->prev = NULL;
    }

    free(temp);
    s->size--;
    return 1;
}

int searchStackByPosition(const Stack* s, size_t position, Citizen* found) {
    Node* node = nodeAtPosition(s, position);

    if (node == NULL) {
        return 0;
    }

    if (found != NULL) {
        *found = node->data;
    }

    return 1;
}

int searchStackBySurname(const Stack* s, const char* surname, Citizen* found, size_t* position) {
    Node* current;
    size_t index;

    if (s == NULL || surname == NULL) {
        return 0;
    }

    current = s->top;
    index = 1;

    while (current != NULL) {
        if (strcmp(current->data.surname, surname) == 0) {
            if (found != NULL) {
                *found = current->data;
            }
            if (position != NULL) {
                *position = index;
            }
            return 1;
        }

        current = current->next;
        index++;
    }

    return 0;
}

int deleteStackByPosition(Stack* s, size_t position, Citizen* removed) {
    Node* target;

    if (s == NULL || position == 0 || position > s->size) {
        return 0;
    }

    if (position == 1) {
        return popStack(s, removed);
    }

    target = nodeAtPosition(s, position);
    if (target == NULL) {
        return 0;
    }

    if (removed != NULL) {
        *removed = target->data;
    }

    if (target->prev != NULL) {
        target->prev->next = target->next;
    }
    if (target->next != NULL) {
        target->next->prev = target->prev;
    }

    free(target);
    s->size--;
    return 1;
}

int deleteStackBySurname(Stack* s, const char* surname, Citizen* removed, size_t* position) {
    Citizen ignored;
    size_t foundPosition;

    if (!searchStackBySurname(s, surname, &ignored, &foundPosition)) {
        return 0;
    }

    if (position != NULL) {
        *position = foundPosition;
    }

    return deleteStackByPosition(s, foundPosition, removed);
}

void displayStack(const Stack* s) {
    Node* current;
    size_t position;

    if (isStackEmpty(s)) {
        printf("Stack is empty.\n");
        return;
    }

    printf("\nVERSION A - Dynamic Stack based on List ADT\n");
    printf("Records: %zu\n", s->size);
    printf("Traversal order: top to bottom\n");
    printf("--------------------------------------------------------------------------\n");

    current = s->top;
    position = 1;
    while (current != NULL) {
        printCitizenBrief(&current->data, position, 0, 0);
        current = current->next;
        position++;
    }

    printf("--------------------------------------------------------------------------\n");
}

void clearStack(Stack* s) {
    if (s == NULL) {
        return;
    }

    while (!isStackEmpty(s)) {
        popStack(s, NULL);
    }
}

\end{lstlisting}

\subsection{queue.h}
\begin{lstlisting}[style=vscodeC]
#ifndef QUEUE_H
#define QUEUE_H

#include "data_types.h"

typedef enum {
    QUEUE_SIMPLE = 1,
    QUEUE_DEQUE = 2,
    QUEUE_CIRCULAR = 3,
    QUEUE_PRIORITY = 4
} QueueType;

typedef struct {
    Node* front;
    Node* rear;
    size_t size;
    QueueType type;
} Queue;

void initQueue(Queue* q, QueueType type);
int isQueueEmpty(const Queue* q);
size_t queueSize(const Queue* q);
const char* queueTypeName(QueueType type);

int enqueueQueue(Queue* q, const Citizen* data, int priority);
int dequeueQueue(Queue* q, Citizen* removed);

int enqueueQueueFront(Queue* q, const Citizen* data);
int dequeueQueueRear(Queue* q, Citizen* removed);

int searchQueueByPosition(const Queue* q, size_t position, Citizen* found, int* priority);
int searchQueueBySurname(const Queue* q, const char* surname, Citizen* found, size_t* position, int* priority);

int deleteQueueByPosition(Queue* q, size_t position, Citizen* removed, int* priority);
int deleteQueueBySurname(Queue* q, const char* surname, Citizen* removed, size_t* position, int* priority);

void displayQueue(const Queue* q);
void clearQueue(Queue* q);

#endif

\end{lstlisting}

\subsection{queue.c}
\begin{lstlisting}[style=vscodeC]
#include <stdio.h>
#include <stdlib.h>
#include <string.h>
#include "queue.h"

static int isCircularQueue(const Queue* q) {
    return q != NULL && q->type == QUEUE_CIRCULAR;
}

static Node* createNode(const Citizen* data, int priority) {
    Node* node = (Node*)malloc(sizeof(Node));
    if (node == NULL) {
        return NULL;
    }

    node->data = *data;
    node->priority = priority;
    node->next = NULL;
    node->prev = NULL;
    return node;
}

static Node* queueNodeAt(const Queue* q, size_t position) {
    Node* current;
    size_t index;

    if (q == NULL || position == 0 || position > q->size) {
        return NULL;
    }

    current = q->front;
    for (index = 1; index < position; index++) {
        current = current->next;
    }

    return current;
}

static int enqueueRearLinear(Queue* q, const Citizen* data, int priority) {
    Node* node = createNode(data, priority);

    if (node == NULL) {
        printf("Memory allocation failed. Queue insert was cancelled.\n");
        return 0;
    }

    node->prev = q->rear;
    if (isQueueEmpty(q)) {
        q->front = node;
    } else {
        q->rear->next = node;
    }

    q->rear = node;
    q->size++;
    return 1;
}

static int enqueueRearCircular(Queue* q, const Citizen* data, int priority) {
    Node* node = createNode(data, priority);

    if (node == NULL) {
        printf("Memory allocation failed. Circular queue insert was cancelled.\n");
        return 0;
    }

    if (isQueueEmpty(q)) {
        q->front = node;
        q->rear = node;
        node->next = node;
        node->prev = node;
    } else {
        node->prev = q->rear;
        node->next = q->front;
        q->rear->next = node;
        q->front->prev = node;
        q->rear = node;
    }

    q->size++;
    return 1;
}

static int enqueueByPriority(Queue* q, const Citizen* data, int priority) {
    Node* node;
    Node* current;

    node = createNode(data, priority);
    if (node == NULL) {
        printf("Memory allocation failed. Priority queue insert was cancelled.\n");
        return 0;
    }

    if (isQueueEmpty(q)) {
        q->front = node;
        q->rear = node;
    } else if (priority < q->front->priority) {
        node->next = q->front;
        q->front->prev = node;
        q->front = node;
    } else {
        current = q->front;
        while (current->next != NULL && current->next->priority <= priority) {
            current = current->next;
        }

        node->next = current->next;
        node->prev = current;
        if (current->next != NULL) {
            current->next->prev = node;
        } else {
            q->rear = node;
        }
        current->next = node;
    }

    q->size++;
    return 1;
}

void initQueue(Queue* q, QueueType type) {
    if (q == NULL) {
        return;
    }

    q->front = NULL;
    q->rear = NULL;
    q->size = 0;
    q->type = type;
}

int isQueueEmpty(const Queue* q) {
    return q == NULL || q->front == NULL;
}

size_t queueSize(const Queue* q) {
    return q == NULL ? 0 : q->size;
}

const char* queueTypeName(QueueType type) {
    switch (type) {
        case QUEUE_SIMPLE:
            return "Simple Queue";
        case QUEUE_DEQUE:
            return "Double Ended Queue";
        case QUEUE_CIRCULAR:
            return "Circular Queue";
        case QUEUE_PRIORITY:
            return "Priority Queue";
        default:
            return "Unknown Queue";
    }
}

int enqueueQueue(Queue* q, const Citizen* data, int priority) {
    if (q == NULL || data == NULL) {
        return 0;
    }

    switch (q->type) {
        case QUEUE_SIMPLE:
        case QUEUE_DEQUE:
            return enqueueRearLinear(q, data, 0);
        case QUEUE_CIRCULAR:
            return enqueueRearCircular(q, data, 0);
        case QUEUE_PRIORITY:
            return enqueueByPriority(q, data, priority);
        default:
            return 0;
    }
}

int enqueueQueueFront(Queue* q, const Citizen* data) {
    Node* node;

    if (q == NULL || data == NULL || q->type != QUEUE_DEQUE) {
        return 0;
    }

    node = createNode(data, 0);
    if (node == NULL) {
        printf("Memory allocation failed. Deque insert was cancelled.\n");
        return 0;
    }

    node->next = q->front;
    if (isQueueEmpty(q)) {
        q->rear = node;
    } else {
        q->front->prev = node;
    }

    q->front = node;
    q->size++;
    return 1;
}

int dequeueQueue(Queue* q, Citizen* removed) {
    Node* temp;

    if (isQueueEmpty(q)) {
        printf("%s underflow. The queue is empty.\n", q == NULL ? "Queue" : queueTypeName(q->type));
        return 0;
    }

    temp = q->front;
    if (removed != NULL) {
        *removed = temp->data;
    }

    if (q->size == 1) {
        q->front = NULL;
        q->rear = NULL;
    } else if (isCircularQueue(q)) {
        q->front = temp->next;
        q->front->prev = q->rear;
        q->rear->next = q->front;
    } else {
        q->front = temp->next;
        q->front->prev = NULL;
    }

    free(temp);
    q->size--;
    return 1;
}

int dequeueQueueRear(Queue* q, Citizen* removed) {
    Node* temp;

    if (q == NULL || q->type != QUEUE_DEQUE) {
        return 0;
    }

    if (isQueueEmpty(q)) {
        printf("Deque underflow. The deque is empty.\n");
        return 0;
    }

    temp = q->rear;
    if (removed != NULL) {
        *removed = temp->data;
    }

    if (q->size == 1) {
        q->front = NULL;
        q->rear = NULL;
    } else {
        q->rear = temp->prev;
        q->rear->next = NULL;
    }

    free(temp);
    q->size--;
    return 1;
}

int searchQueueByPosition(const Queue* q, size_t position, Citizen* found, int* priority) {
    Node* node = queueNodeAt(q, position);

    if (node == NULL) {
        return 0;
    }

    if (found != NULL) {
        *found = node->data;
    }
    if (priority != NULL) {
        *priority = node->priority;
    }

    return 1;
}

int searchQueueBySurname(const Queue* q, const char* surname, Citizen* found, size_t* position, int* priority) {
    Node* current;
    size_t index;

    if (q == NULL || surname == NULL) {
        return 0;
    }

    current = q->front;
    for (index = 1; index <= q->size; index++) {
        if (strcmp(current->data.surname, surname) == 0) {
            if (found != NULL) {
                *found = current->data;
            }
            if (position != NULL) {
                *position = index;
            }
            if (priority != NULL) {
                *priority = current->priority;
            }
            return 1;
        }

        current = current->next;
    }

    return 0;
}

int deleteQueueByPosition(Queue* q, size_t position, Citizen* removed, int* priority) {
    Node* target;

    if (q == NULL || position == 0 || position > q->size) {
        return 0;
    }

    if (position == 1) {
        target = q->front;
        if (priority != NULL) {
            *priority = target->priority;
        }
        return dequeueQueue(q, removed);
    }

    target = queueNodeAt(q, position);
    if (target == NULL) {
        return 0;
    }

    if (removed != NULL) {
        *removed = target->data;
    }
    if (priority != NULL) {
        *priority = target->priority;
    }

    target->prev->next = target->next;
    if (isCircularQueue(q)) {
        target->next->prev = target->prev;
        if (target == q->rear) {
            q->rear = target->prev;
        }
    } else if (target->next != NULL) {
        target->next->prev = target->prev;
    } else {
        q->rear = target->prev;
    }

    free(target);
    q->size--;
    return 1;
}

int deleteQueueBySurname(Queue* q, const char* surname, Citizen* removed, size_t* position, int* priority) {
    Citizen ignored;
    size_t foundPosition;

    if (!searchQueueBySurname(q, surname, &ignored, &foundPosition, priority)) {
        return 0;
    }

    if (position != NULL) {
        *position = foundPosition;
    }

    return deleteQueueByPosition(q, foundPosition, removed, priority);
}

void displayQueue(const Queue* q) {
    Node* current;
    size_t position;
    int showPriority;

    if (isQueueEmpty(q)) {
        printf("%s is empty.\n", q == NULL ? "Queue" : queueTypeName(q->type));
        return;
    }

    showPriority = q->type == QUEUE_PRIORITY;
    printf("\nVERSION B - %s based on List ADT\n", queueTypeName(q->type));
    printf("Records: %zu\n", q->size);
    printf("Traversal order: front to rear\n");
    printf("--------------------------------------------------------------------------\n");

    current = q->front;
    for (position = 1; position <= q->size; position++) {
        printCitizenBrief(&current->data, position, current->priority, showPriority);
        current = current->next;
    }

    printf("--------------------------------------------------------------------------\n");
}

void clearQueue(Queue* q) {
    if (q == NULL) {
        return;
    }

    while (!isQueueEmpty(q)) {
        dequeueQueue(q, NULL);
    }
}

\end{lstlisting}

\subsection{file\_io.h}
\begin{lstlisting}[style=vscodeC]
#ifndef FILE_IO_H
#define FILE_IO_H

#include "stack.h"
#include "queue.h"

int createEmptyDataFile(const char* filename, int binaryMode);

int saveStackToTextFile(const Stack* s, const char* filename);
int saveStackToBinaryFile(const Stack* s, const char* filename);
int loadStackFromBinaryFile(Stack* s, const char* filename, int replaceExisting);

int saveQueueToTextFile(const Queue* q, const char* filename);
int saveQueueToBinaryFile(const Queue* q, const char* filename);
int loadQueueFromBinaryFile(Queue* q, const char* filename, int replaceExisting);

int displayTextFile(const char* filename);
int displayStackBinaryFile(const char* filename);
int displayQueueBinaryFile(const char* filename);
int deleteDataFile(const char* filename);

#endif

\end{lstlisting}

\subsection{file\_io.c}
\begin{lstlisting}[style=vscodeC]
#include <stdio.h>
#include <stdlib.h>
#include "file_io.h"

typedef struct {
    Citizen data;
    int priority;
} QueueFileRecord;

static void writeCitizenText(FILE* file, const Citizen* citizen, int priority, int showPriority) {
    fprintf(file,
            "%s %s | DOB: %02d/%02d/%04d | Gender: %c | Age: %dY %dM %dD | Category: %s | Paid: %.2f",
            citizen->name,
            citizen->surname,
            citizen->dob.day,
            citizen->dob.month,
            citizen->dob.year,
            citizen->gender,
            citizen->ageYears,
            citizen->ageMonths,
            citizen->ageDays,
            citizenCategoryName(citizen->category),
            citizen->amountPaid);

    if (showPriority) {
        fprintf(file, " | Priority: %d", priority);
    }

    fprintf(file,
            " | Home: %s, %s, %s | Work: %s, %s, %s\n",
            citizen->home.city,
            citizen->home.street,
            citizen->home.postCode,
            citizen->work.city,
            citizen->work.street,
            citizen->work.postCode);
}

int createEmptyDataFile(const char* filename, int binaryMode) {
    FILE* file;

    if (filename == NULL) {
        return 0;
    }

    file = fopen(filename, binaryMode ? "wb" : "w");
    if (file == NULL) {
        printf("Could not create file: %s\n", filename);
        return 0;
    }

    fclose(file);
    printf("File created: %s\n", filename);
    return 1;
}

int saveStackToTextFile(const Stack* s, const char* filename) {
    FILE* file;
    Node* current;
    size_t position;

    if (s == NULL || filename == NULL) {
        return 0;
    }

    file = fopen(filename, "w");
    if (file == NULL) {
        printf("Could not open file for text write: %s\n", filename);
        return 0;
    }

    fprintf(file, "VERSION A - Dynamic Stack based on List ADT\n");
    fprintf(file, "Records: %zu\n\n", s->size);

    current = s->top;
    position = 1;
    while (current != NULL) {
        fprintf(file, "%zu. ", position);
        writeCitizenText(file, &current->data, 0, 0);
        current = current->next;
        position++;
    }

    fclose(file);
    printf("Stack saved in text mode: %s\n", filename);
    return 1;
}

int saveStackToBinaryFile(const Stack* s, const char* filename) {
    FILE* file;
    Node* current;

    if (s == NULL || filename == NULL) {
        return 0;
    }

    file = fopen(filename, "wb");
    if (file == NULL) {
        printf("Could not open file for binary write: %s\n", filename);
        return 0;
    }

    current = s->top;
    while (current != NULL) {
        fwrite(&current->data, sizeof(Citizen), 1, file);
        current = current->next;
    }

    fclose(file);
    printf("Stack saved in binary mode: %s\n", filename);
    return 1;
}

int loadStackFromBinaryFile(Stack* s, const char* filename, int replaceExisting) {
    FILE* file;
    Citizen* records;
    Citizen* resized;
    Citizen item;
    size_t count;
    size_t capacity;
    size_t index;

    if (s == NULL || filename == NULL) {
        return 0;
    }

    file = fopen(filename, "rb");
    if (file == NULL) {
        printf("Could not open binary stack file: %s\n", filename);
        return 0;
    }

    capacity = 8;
    count = 0;
    records = (Citizen*)malloc(capacity * sizeof(Citizen));
    if (records == NULL) {
        fclose(file);
        printf("Memory allocation failed while loading stack file.\n");
        return 0;
    }

    while (fread(&item, sizeof(Citizen), 1, file) == 1) {
        if (count == capacity) {
            capacity *= 2;
            resized = (Citizen*)realloc(records, capacity * sizeof(Citizen));
            if (resized == NULL) {
                free(records);
                fclose(file);
                printf("Memory allocation failed while expanding stack load buffer.\n");
                return 0;
            }
            records = resized;
        }

        records[count++] = item;
    }

    fclose(file);

    if (replaceExisting) {
        clearStack(s);
    }

    for (index = count; index > 0; index--) {
        if (!pushStack(s, &records[index - 1])) {
            free(records);
            printf("Stack load stopped because an insert failed.\n");
            return 0;
        }
    }

    free(records);
    printf("Loaded %zu stack records from: %s\n", count, filename);
    return 1;
}

int saveQueueToTextFile(const Queue* q, const char* filename) {
    FILE* file;
    Node* current;
    size_t position;
    int showPriority;

    if (q == NULL || filename == NULL) {
        return 0;
    }

    file = fopen(filename, "w");
    if (file == NULL) {
        printf("Could not open file for text write: %s\n", filename);
        return 0;
    }

    showPriority = q->type == QUEUE_PRIORITY;
    fprintf(file, "VERSION B - %s based on List ADT\n", queueTypeName(q->type));
    fprintf(file, "Records: %zu\n\n", q->size);

    current = q->front;
    for (position = 1; position <= q->size; position++) {
        fprintf(file, "%zu. ", position);
        writeCitizenText(file, &current->data, current->priority, showPriority);
        current = current->next;
    }

    fclose(file);
    printf("Queue saved in text mode: %s\n", filename);
    return 1;
}

int saveQueueToBinaryFile(const Queue* q, const char* filename) {
    FILE* file;
    Node* current;
    QueueFileRecord record;
    size_t position;

    if (q == NULL || filename == NULL) {
        return 0;
    }

    file = fopen(filename, "wb");
    if (file == NULL) {
        printf("Could not open file for binary write: %s\n", filename);
        return 0;
    }

    current = q->front;
    for (position = 1; position <= q->size; position++) {
        record.data = current->data;
        record.priority = current->priority;
        fwrite(&record, sizeof(QueueFileRecord), 1, file);
        current = current->next;
    }

    fclose(file);
    printf("Queue saved in binary mode: %s\n", filename);
    return 1;
}

int loadQueueFromBinaryFile(Queue* q, const char* filename, int replaceExisting) {
    FILE* file;
    QueueFileRecord record;
    size_t count;

    if (q == NULL || filename == NULL) {
        return 0;
    }

    file = fopen(filename, "rb");
    if (file == NULL) {
        printf("Could not open binary queue file: %s\n", filename);
        return 0;
    }

    if (replaceExisting) {
        clearQueue(q);
    }

    count = 0;
    while (fread(&record, sizeof(QueueFileRecord), 1, file) == 1) {
        if (!enqueueQueue(q, &record.data, record.priority)) {
            fclose(file);
            printf("Queue load stopped because an insert failed.\n");
            return 0;
        }
        count++;
    }

    fclose(file);
    printf("Loaded %zu queue records from: %s\n", count, filename);
    return 1;
}

int displayTextFile(const char* filename) {
    FILE* file;
    int ch;

    if (filename == NULL) {
        return 0;
    }

    file = fopen(filename, "r");
    if (file == NULL) {
        printf("Could not open text file: %s\n", filename);
        return 0;
    }

    printf("\n--- Text file: %s ---\n", filename);
    while ((ch = fgetc(file)) != EOF) {
        putchar(ch);
    }
    printf("\n--- End of text file ---\n");

    fclose(file);
    return 1;
}

int displayStackBinaryFile(const char* filename) {
    FILE* file;
    Citizen item;
    size_t position;

    if (filename == NULL) {
        return 0;
    }

    file = fopen(filename, "rb");
    if (file == NULL) {
        printf("Could not open binary stack file: %s\n", filename);
        return 0;
    }

    printf("\n--- Binary stack file: %s ---\n", filename);
    position = 1;
    while (fread(&item, sizeof(Citizen), 1, file) == 1) {
        printCitizenBrief(&item, position, 0, 0);
        position++;
    }
    printf("--- End of binary stack file ---\n");

    fclose(file);
    return 1;
}

int displayQueueBinaryFile(const char* filename) {
    FILE* file;
    QueueFileRecord record;
    size_t position;

    if (filename == NULL) {
        return 0;
    }

    file = fopen(filename, "rb");
    if (file == NULL) {
        printf("Could not open binary queue file: %s\n", filename);
        return 0;
    }

    printf("\n--- Binary queue file: %s ---\n", filename);
    position = 1;
    while (fread(&record, sizeof(QueueFileRecord), 1, file) == 1) {
        printCitizenBrief(&record.data, position, record.priority, 1);
        position++;
    }
    printf("--- End of binary queue file ---\n");

    fclose(file);
    return 1;
}

int deleteDataFile(const char* filename) {
    if (filename == NULL) {
        return 0;
    }

    if (remove(filename) == 0) {
        printf("File deleted: %s\n", filename);
        return 1;
    }

    printf("Could not delete file: %s\n", filename);
    return 0;
}

\end{lstlisting}

\subsection{main.c}
\begin{lstlisting}[style=vscodeC]
#include <ctype.h>
#include <errno.h>
#include <stdio.h>
#include <stdlib.h>
#include <string.h>
#include <time.h>
#include "data_types.h"
#include "file_io.h"
#include "queue.h"
#include "stack.h"

#define INPUT_LENGTH 256
#define MAX_RECORDS 10000

static void trimNewline(char* text) {
    if (text != NULL) {
        text[strcspn(text, "\n")] = '\0';
    }
}

static int readLine(const char* prompt, char* buffer, size_t length) {
    if (prompt != NULL) {
        printf("%s", prompt);
    }

    if (fgets(buffer, (int)length, stdin) == NULL) {
        return 0;
    }

    trimNewline(buffer);
    return 1;
}

static void readRequiredString(const char* prompt, char* destination, size_t length) {
    char buffer[INPUT_LENGTH];

    do {
        readLine(prompt, buffer, sizeof(buffer));
        if (buffer[0] == '\0') {
            printf("This field cannot be empty.\n");
        }
    } while (buffer[0] == '\0');

    strncpy(destination, buffer, length - 1);
    destination[length - 1] = '\0';
}

static int readIntRange(const char* prompt, int minValue, int maxValue) {
    char buffer[INPUT_LENGTH];
    char* end;
    long value;

    while (1) {
        readLine(prompt, buffer, sizeof(buffer));
        errno = 0;
        value = strtol(buffer, &end, 10);
        while (isspace((unsigned char)*end)) {
            end++;
        }

        if (errno == 0 && end != buffer && *end == '\0' && value >= minValue && value <= maxValue) {
            return (int)value;
        }

        printf("Enter an integer from %d to %d.\n", minValue, maxValue);
    }
}

static float readFloatMin(const char* prompt, float minValue) {
    char buffer[INPUT_LENGTH];
    char* end;
    float value;

    while (1) {
        readLine(prompt, buffer, sizeof(buffer));
        errno = 0;
        value = strtof(buffer, &end);
        while (isspace((unsigned char)*end)) {
            end++;
        }

        if (errno == 0 && end != buffer && *end == '\0' && value >= minValue) {
            return value;
        }

        printf("Enter a number greater than or equal to %.2f.\n", minValue);
    }
}

static int readYesNo(const char* prompt) {
    char buffer[INPUT_LENGTH];

    while (1) {
        readLine(prompt, buffer, sizeof(buffer));
        if (buffer[0] == 'y' || buffer[0] == 'Y') {
            return 1;
        }
        if (buffer[0] == 'n' || buffer[0] == 'N') {
            return 0;
        }

        printf("Enter y or n.\n");
    }
}

static Date currentDate(void) {
    time_t rawTime = time(NULL);
    struct tm* local = localtime(&rawTime);
    Date today = {8, 5, 2026};

    if (local != NULL) {
        today.day = local->tm_mday;
        today.month = local->tm_mon + 1;
        today.year = local->tm_year + 1900;
    }

    return today;
}

static int isLeapYear(int year) {
    return (year % 400 == 0) || (year % 4 == 0 && year % 100 != 0);
}

static int daysInMonth(int month, int year) {
    static const int days[] = {31, 28, 31, 30, 31, 30, 31, 31, 30, 31, 30, 31};

    if (month == 2 && isLeapYear(year)) {
        return 29;
    }

    return days[month - 1];
}

static int isFutureDate(Date date, Date today) {
    if (date.year != today.year) {
        return date.year > today.year;
    }
    if (date.month != today.month) {
        return date.month > today.month;
    }
    return date.day > today.day;
}

static int isValidDate(Date date) {
    if (date.year < 1900 || date.month < 1 || date.month > 12 || date.day < 1) {
        return 0;
    }

    return date.day <= daysInMonth(date.month, date.year);
}

static Date readDateOfBirth(void) {
    Date date;
    Date today = currentDate();

    while (1) {
        date.day = readIntRange("Birth day: ", 1, 31);
        date.month = readIntRange("Birth month: ", 1, 12);
        date.year = readIntRange("Birth year: ", 1900, today.year);

        if (isValidDate(date) && !isFutureDate(date, today)) {
            return date;
        }

        printf("Invalid or future date. Please enter the date again.\n");
    }
}

static void updateDerivedCitizenFields(Citizen* citizen) {
    Date today;
    int previousMonth;
    int previousYear;

    if (citizen == NULL) {
        return;
    }

    today = currentDate();
    citizen->ageYears = today.year - citizen->dob.year;
    citizen->ageMonths = today.month - citizen->dob.month;
    citizen->ageDays = today.day - citizen->dob.day;

    if (citizen->ageDays < 0) {
        citizen->ageMonths--;
        previousMonth = today.month - 1;
        previousYear = today.year;
        if (previousMonth == 0) {
            previousMonth = 12;
            previousYear--;
        }
        citizen->ageDays += daysInMonth(previousMonth, previousYear);
    }

    if (citizen->ageMonths < 0) {
        citizen->ageYears--;
        citizen->ageMonths += 12;
    }

    if (citizen->ageYears < 18) {
        citizen->category = CATEGORY_CHILD;
    } else if (citizen->ageYears < 60) {
        citizen->category = CATEGORY_ADULT;
    } else {
        citizen->category = CATEGORY_SENIOR;
    }
}

static char readGender(void) {
    char buffer[INPUT_LENGTH];
    char gender;

    while (1) {
        readLine("Gender (M/F/O): ", buffer, sizeof(buffer));
        gender = (char)toupper((unsigned char)buffer[0]);
        if (gender == 'M' || gender == 'F' || gender == 'O') {
            return gender;
        }

        printf("Enter M, F, or O.\n");
    }
}

static void inputAddress(const char* title, Address* address) {
    printf("%s\n", title);
    readRequiredString("  City: ", address->city, CITY_LENGTH);
    readRequiredString("  Street: ", address->street, STREET_LENGTH);
    readRequiredString("  Postal code: ", address->postCode, POST_CODE_LENGTH);
}

static void inputCitizen(Citizen* citizen, size_t ordinal) {
    if (citizen == NULL) {
        return;
    }

    memset(citizen, 0, sizeof(Citizen));
    printf("\nCitizen %zu\n", ordinal);
    readRequiredString("Name: ", citizen->name, NAME_LENGTH);
    readRequiredString("Surname: ", citizen->surname, SURNAME_LENGTH);
    citizen->dob = readDateOfBirth();
    citizen->gender = readGender();
    inputAddress("Home address", &citizen->home);
    inputAddress("Work address", &citizen->work);
    citizen->amountPaid = readFloatMin("Amount paid: ", 0.0f);
    updateDerivedCitizenFields(citizen);
}

static void readFileAddress(char* filename, size_t length) {
    readRequiredString("File address/name: ", filename, length);
}

static void registerStackToFile(const Stack* stack) {
    char filename[INPUT_LENGTH];
    int mode;

    mode = readIntRange("Register stack as 1.Text or 2.Binary: ", 1, 2);
    readFileAddress(filename, sizeof(filename));

    if (mode == 1) {
        saveStackToTextFile(stack, filename);
    } else {
        saveStackToBinaryFile(stack, filename);
    }
}

static void registerQueueToFile(const Queue* queue) {
    char filename[INPUT_LENGTH];
    int mode;

    mode = readIntRange("Register queue as 1.Text or 2.Binary: ", 1, 2);
    readFileAddress(filename, sizeof(filename));

    if (mode == 1) {
        saveQueueToTextFile(queue, filename);
    } else {
        saveQueueToBinaryFile(queue, filename);
    }
}

static void offerStackRegistration(const Stack* stack) {
    if (readYesNo("Register the current stack in a file now? (y/n): ")) {
        registerStackToFile(stack);
    }
}

static void offerQueueRegistration(const Queue* queue) {
    if (readYesNo("Register the current queue in a file now? (y/n): ")) {
        registerQueueToFile(queue);
    }
}

static void createStackFromRecords(Stack* stack) {
    Citizen* records;
    size_t count;
    size_t index;

    count = (size_t)readIntRange("Number of stack records: ", 1, MAX_RECORDS);
    records = (Citizen*)malloc(count * sizeof(Citizen));
    if (records == NULL) {
        printf("Memory allocation failed. Stack creation was cancelled.\n");
        return;
    }

    for (index = 0; index < count; index++) {
        inputCitizen(&records[index], index + 1);
    }

    clearStack(stack);
    for (index = 0; index < count; index++) {
        pushStack(stack, &records[index]);
    }

    free(records);
    printf("Stack created with %zu records. The last entered record is on top.\n", stackSize(stack));
    offerStackRegistration(stack);
}

static void pushOneCitizen(Stack* stack) {
    Citizen citizen;

    inputCitizen(&citizen, stackSize(stack) + 1);
    if (pushStack(stack, &citizen)) {
        printf("Citizen pushed to the stack.\n");
        offerStackRegistration(stack);
    }
}

static void popOneCitizen(Stack* stack) {
    Citizen removed;

    if (popStack(stack, &removed)) {
        printf("Removed from stack top:\n");
        printCitizenDetails(&removed);
        offerStackRegistration(stack);
    }
}

static void searchStackPositionMenu(const Stack* stack) {
    Citizen found;
    int position;

    if (isStackEmpty(stack)) {
        printf("Stack is empty.\n");
        return;
    }

    position = readIntRange("Position from top: ", 1, (int)stackSize(stack));
    if (searchStackByPosition(stack, (size_t)position, &found)) {
        printCitizenDetails(&found);
    } else {
        printf("No stack record found at that position.\n");
    }
}

static void searchStackSurnameMenu(const Stack* stack) {
    char surname[SURNAME_LENGTH];
    Citizen found;
    size_t position;

    readRequiredString("Surname to search: ", surname, sizeof(surname));
    if (searchStackBySurname(stack, surname, &found, &position)) {
        printf("Found at position %zu from top.\n", position);
        printCitizenDetails(&found);
    } else {
        printf("No stack record found for surname: %s\n", surname);
    }
}

static void deleteStackPositionMenu(Stack* stack) {
    Citizen removed;
    int position;

    if (isStackEmpty(stack)) {
        printf("Stack is empty.\n");
        return;
    }

    position = readIntRange("Position to delete from top: ", 1, (int)stackSize(stack));
    if (deleteStackByPosition(stack, (size_t)position, &removed)) {
        printf("Deleted stack record:\n");
        printCitizenDetails(&removed);
        offerStackRegistration(stack);
    } else {
        printf("Delete failed.\n");
    }
}

static void deleteStackSurnameMenu(Stack* stack) {
    char surname[SURNAME_LENGTH];
    Citizen removed;
    size_t position;

    readRequiredString("Surname to delete: ", surname, sizeof(surname));
    if (deleteStackBySurname(stack, surname, &removed, &position)) {
        printf("Deleted record at position %zu from top:\n", position);
        printCitizenDetails(&removed);
        offerStackRegistration(stack);
    } else {
        printf("No stack record found for surname: %s\n", surname);
    }
}

static void loadStackMenu(Stack* stack) {
    char filename[INPUT_LENGTH];
    int replace;

    readFileAddress(filename, sizeof(filename));
    replace = readYesNo("Replace the current stack before loading? (y/n): ");
    loadStackFromBinaryFile(stack, filename, replace);
}

static void versionAStackMenu(Stack* stack) {
    int choice;
    char filename[INPUT_LENGTH];

    while (1) {
        printf("\n================ VERSION A =================\n");
        printf("Dynamic Stack based on List ADT\n");
        printf("1. Create or replace stack from N records\n");
        printf("2. Insert element: push to stack\n");
        printf("3. Delete element: pop from stack\n");
        printf("4. Traverse and display stack\n");
        printf("5. Search by position\n");
        printf("6. Search by surname\n");
        printf("7. Delete by position\n");
        printf("8. Delete by surname\n");
        printf("9. Save stack to text file\n");
        printf("10. Save stack to binary file\n");
        printf("11. Load stack from binary file\n");
        printf("0. Back to main menu\n");

        choice = readIntRange("Choice: ", 0, 11);
        switch (choice) {
            case 1:
                createStackFromRecords(stack);
                break;
            case 2:
                pushOneCitizen(stack);
                break;
            case 3:
                popOneCitizen(stack);
                break;
            case 4:
                displayStack(stack);
                break;
            case 5:
                searchStackPositionMenu(stack);
                break;
            case 6:
                searchStackSurnameMenu(stack);
                break;
            case 7:
                deleteStackPositionMenu(stack);
                break;
            case 8:
                deleteStackSurnameMenu(stack);
                break;
            case 9:
                readFileAddress(filename, sizeof(filename));
                saveStackToTextFile(stack, filename);
                break;
            case 10:
                readFileAddress(filename, sizeof(filename));
                saveStackToBinaryFile(stack, filename);
                break;
            case 11:
                loadStackMenu(stack);
                break;
            case 0:
                return;
            default:
                break;
        }
    }
}

static QueueType readQueueType(void) {
    printf("\nQueue type\n");
    printf("1. Simple Queue\n");
    printf("2. Double Ended Queue\n");
    printf("3. Circular Queue\n");
    printf("4. Priority Queue\n");
    return (QueueType)readIntRange("Select queue type: ", 1, 4);
}

static Queue* queueByType(Queue queues[], QueueType type) {
    return &queues[(int)type - 1];
}

static int readPriority(void) {
    return readIntRange("Priority (1 = highest priority): ", 1, 1000000);
}

static void createQueueFromRecords(Queue* queue) {
    Citizen* records;
    int* priorities;
    size_t count;
    size_t index;

    count = (size_t)readIntRange("Number of queue records: ", 1, MAX_RECORDS);
    records = (Citizen*)malloc(count * sizeof(Citizen));
    priorities = (int*)malloc(count * sizeof(int));

    if (records == NULL || priorities == NULL) {
        free(records);
        free(priorities);
        printf("Memory allocation failed. Queue creation was cancelled.\n");
        return;
    }

    for (index = 0; index < count; index++) {
        inputCitizen(&records[index], index + 1);
        priorities[index] = queue->type == QUEUE_PRIORITY ? readPriority() : 0;
    }

    clearQueue(queue);
    for (index = 0; index < count; index++) {
        enqueueQueue(queue, &records[index], priorities[index]);
    }

    free(records);
    free(priorities);
    printf("%s created with %zu records.\n", queueTypeName(queue->type), queueSize(queue));
    offerQueueRegistration(queue);
}

static void enqueueOneCitizen(Queue* queue) {
    Citizen citizen;
    int option;
    int priority;

    inputCitizen(&citizen, queueSize(queue) + 1);

    if (queue->type == QUEUE_DEQUE) {
        option = readIntRange("Insert at 1.Front or 2.Rear: ", 1, 2);
        if (option == 1) {
            if (enqueueQueueFront(queue, &citizen)) {
                printf("Citizen inserted at the deque front.\n");
                offerQueueRegistration(queue);
            }
        } else if (enqueueQueue(queue, &citizen, 0)) {
            printf("Citizen inserted at the deque rear.\n");
            offerQueueRegistration(queue);
        }
        return;
    }

    priority = queue->type == QUEUE_PRIORITY ? readPriority() : 0;
    if (enqueueQueue(queue, &citizen, priority)) {
        printf("Citizen inserted into %s.\n", queueTypeName(queue->type));
        offerQueueRegistration(queue);
    }
}

static void dequeueOneCitizen(Queue* queue) {
    Citizen removed;
    int option;

    if (queue->type == QUEUE_DEQUE) {
        option = readIntRange("Delete from 1.Front or 2.Rear: ", 1, 2);
        if (option == 2) {
            if (dequeueQueueRear(queue, &removed)) {
                printf("Removed from deque rear:\n");
                printCitizenDetails(&removed);
                offerQueueRegistration(queue);
            }
            return;
        }
    }

    if (dequeueQueue(queue, &removed)) {
        printf("Removed from queue front:\n");
        printCitizenDetails(&removed);
        offerQueueRegistration(queue);
    }
}

static void searchQueuePositionMenu(const Queue* queue) {
    Citizen found;
    int priority;
    int position;

    if (isQueueEmpty(queue)) {
        printf("%s is empty.\n", queueTypeName(queue->type));
        return;
    }

    position = readIntRange("Position from front: ", 1, (int)queueSize(queue));
    if (searchQueueByPosition(queue, (size_t)position, &found, &priority)) {
        if (queue->type == QUEUE_PRIORITY) {
            printf("Priority: %d\n", priority);
        }
        printCitizenDetails(&found);
    } else {
        printf("No queue record found at that position.\n");
    }
}

static void searchQueueSurnameMenu(const Queue* queue) {
    char surname[SURNAME_LENGTH];
    Citizen found;
    size_t position;
    int priority;

    readRequiredString("Surname to search: ", surname, sizeof(surname));
    if (searchQueueBySurname(queue, surname, &found, &position, &priority)) {
        printf("Found at position %zu from front.\n", position);
        if (queue->type == QUEUE_PRIORITY) {
            printf("Priority: %d\n", priority);
        }
        printCitizenDetails(&found);
    } else {
        printf("No queue record found for surname: %s\n", surname);
    }
}

static void deleteQueuePositionMenu(Queue* queue) {
    Citizen removed;
    int priority;
    int position;

    if (isQueueEmpty(queue)) {
        printf("%s is empty.\n", queueTypeName(queue->type));
        return;
    }

    position = readIntRange("Position to delete from front: ", 1, (int)queueSize(queue));
    if (deleteQueueByPosition(queue, (size_t)position, &removed, &priority)) {
        printf("Deleted queue record:\n");
        if (queue->type == QUEUE_PRIORITY) {
            printf("Priority: %d\n", priority);
        }
        printCitizenDetails(&removed);
        offerQueueRegistration(queue);
    } else {
        printf("Delete failed.\n");
    }
}

static void deleteQueueSurnameMenu(Queue* queue) {
    char surname[SURNAME_LENGTH];
    Citizen removed;
    size_t position;
    int priority;

    readRequiredString("Surname to delete: ", surname, sizeof(surname));
    if (deleteQueueBySurname(queue, surname, &removed, &position, &priority)) {
        printf("Deleted record at position %zu from front:\n", position);
        if (queue->type == QUEUE_PRIORITY) {
            printf("Priority: %d\n", priority);
        }
        printCitizenDetails(&removed);
        offerQueueRegistration(queue);
    } else {
        printf("No queue record found for surname: %s\n", surname);
    }
}

static void loadQueueMenu(Queue* queue) {
    char filename[INPUT_LENGTH];
    int replace;

    readFileAddress(filename, sizeof(filename));
    replace = readYesNo("Replace the current queue before loading? (y/n): ");
    loadQueueFromBinaryFile(queue, filename, replace);
}

static void versionBQueueMenu(Queue queues[]) {
    QueueType activeType = readQueueType();
    Queue* activeQueue = queueByType(queues, activeType);
    int choice;
    char filename[INPUT_LENGTH];

    while (1) {
        printf("\n================ VERSION B =================\n");
        printf("%s based on List ADT\n", queueTypeName(activeType));
        printf("1. Switch queue type\n");
        printf("2. Create or replace selected queue from N records\n");
        printf("3. Insert element into selected queue\n");
        printf("4. Delete element from selected queue\n");
        printf("5. Traverse and display selected queue\n");
        printf("6. Search by position\n");
        printf("7. Search by surname\n");
        printf("8. Delete by position\n");
        printf("9. Delete by surname\n");
        printf("10. Save selected queue to text file\n");
        printf("11. Save selected queue to binary file\n");
        printf("12. Load selected queue from binary file\n");
        printf("0. Back to main menu\n");

        choice = readIntRange("Choice: ", 0, 12);
        switch (choice) {
            case 1:
                activeType = readQueueType();
                activeQueue = queueByType(queues, activeType);
                break;
            case 2:
                createQueueFromRecords(activeQueue);
                break;
            case 3:
                enqueueOneCitizen(activeQueue);
                break;
            case 4:
                dequeueOneCitizen(activeQueue);
                break;
            case 5:
                displayQueue(activeQueue);
                break;
            case 6:
                searchQueuePositionMenu(activeQueue);
                break;
            case 7:
                searchQueueSurnameMenu(activeQueue);
                break;
            case 8:
                deleteQueuePositionMenu(activeQueue);
                break;
            case 9:
                deleteQueueSurnameMenu(activeQueue);
                break;
            case 10:
                readFileAddress(filename, sizeof(filename));
                saveQueueToTextFile(activeQueue, filename);
                break;
            case 11:
                readFileAddress(filename, sizeof(filename));
                saveQueueToBinaryFile(activeQueue, filename);
                break;
            case 12:
                loadQueueMenu(activeQueue);
                break;
            case 0:
                return;
            default:
                break;
        }
    }
}

static void fileToolsMenu(void) {
    int choice;
    char filename[INPUT_LENGTH];

    while (1) {
        printf("\n================ FILE TOOLS =================\n");
        printf("1. Create empty text file\n");
        printf("2. Create empty binary file\n");
        printf("3. Open or reopen text file\n");
        printf("4. Open or reopen binary stack file\n");
        printf("5. Open or reopen binary queue file\n");
        printf("6. Delete file\n");
        printf("0. Back to main menu\n");

        choice = readIntRange("Choice: ", 0, 6);
        switch (choice) {
            case 1:
                readFileAddress(filename, sizeof(filename));
                createEmptyDataFile(filename, 0);
                break;
            case 2:
                readFileAddress(filename, sizeof(filename));
                createEmptyDataFile(filename, 1);
                break;
            case 3:
                readFileAddress(filename, sizeof(filename));
                displayTextFile(filename);
                break;
            case 4:
                readFileAddress(filename, sizeof(filename));
                displayStackBinaryFile(filename);
                break;
            case 5:
                readFileAddress(filename, sizeof(filename));
                displayQueueBinaryFile(filename);
                break;
            case 6:
                readFileAddress(filename, sizeof(filename));
                deleteDataFile(filename);
                break;
            case 0:
                return;
            default:
                break;
        }
    }
}

int main(void) {
    Stack stackVersionA;
    Queue queuesVersionB[4];
    int choice;
    size_t index;

    initStack(&stackVersionA);
    initQueue(&queuesVersionB[0], QUEUE_SIMPLE);
    initQueue(&queuesVersionB[1], QUEUE_DEQUE);
    initQueue(&queuesVersionB[2], QUEUE_CIRCULAR);
    initQueue(&queuesVersionB[3], QUEUE_PRIORITY);

    while (1) {
        printf("\n================ MAIN MENU =================\n");
        printf("1. Version A: Dynamic Stack based on List ADT\n");
        printf("2. Version B: Dynamic Queues based on List ADT\n");
        printf("3. File create/open/reopen/delete tools\n");
        printf("0. Exit\n");

        choice = readIntRange("Choice: ", 0, 3);
        switch (choice) {
            case 1:
                versionAStackMenu(&stackVersionA);
                break;
            case 2:
                versionBQueueMenu(queuesVersionB);
                break;
            case 3:
                fileToolsMenu();
                break;
            case 0:
                clearStack(&stackVersionA);
                for (index = 0; index < 4; index++) {
                    clearQueue(&queuesVersionB[index]);
                }
                printf("Memory released. Program finished.\n");
                return 0;
            default:
                break;
        }
    }
}

\end{lstlisting}

\clearpage

\section{Conclusion}

This laboratory work presents a complete procedural C solution for processing user-defined records with Stack
and Queue abstract data types. The program is based on the \texttt{Citizen} custom data type, which contains
initial input fields and computed output fields. The record also contains nested structures for date and
addresses, which makes the model closer to a real data-processing task.

Version A implements a dynamic Stack based on a linked-list ADT. It supports creation from a user-defined
number of records, insertion through push, deletion through pop or direct removal, traversal, searching by
position, searching by surname, and file registration in text or binary mode.

Version B implements four dynamic Queues based on a linked-list ADT: Simple Queue, Double Ended Queue,
Circular Queue, and Priority Queue. The same menu structure allows the user to choose the queue type and then
perform creation, insertion, deletion, traversal, searching, and file registration operations. The Priority
Queue additionally stores a priority value and keeps records ordered so that the highest-priority citizen is
processed first.

The solution uses pointers throughout the implementation. List nodes are dynamically allocated with
\texttt{malloc}, resized loading buffers use \texttt{realloc}, and memory is released with \texttt{free}. File
processing is implemented through the C \texttt{FILE} data type, supporting both text and binary output, binary
loading, text-file viewing, binary-file viewing, file creation, file reopening, and file deletion.

Overall, the final program satisfies the stated requirements: it is procedural, modular, menu-driven, organized
through header files, based on user-defined structures, uses dynamic memory and pointers, implements the
required Stack and Queue variants, and includes the requested file operations.

}
\end{document}

